Para a análise da solução desenvolvida e identificação de problemas foram feitos testes variados. Este capítulo foi dividido em seções que representam as descobertas mais relevantes ao longo da pesquisa e desenvolvimento. A seção~\ref{sec:extracao_da_placa_resultados} vai descrever os resultados obtidos na experimentação com a extração da placa nas imagens, a seção~\ref{sec:reconhecimento_dos_caracteres_resultados} vai abordar os resultados obtidos no reconhecimento dos caracteres e a seção~\ref{sec:performance_resultados} vai abordar como performa o algoritmo no computador \emph{Raspberry Pi}.

Foram coletadas imagens de N veículos distintos para analisar a capacidade de reconhecimento do sistema. A tabela~\ref{tab:resultados} mostra o número de imagens processadas, o número de imagens onde a placa foi encontrada com sucesso e o número de imagens onde foi possível reconhecer o valor da placa.

\begin{table}[]
\centering
\caption{Resultados obtidos}
\label{tab:resultados}
\begin{tabular}{@{}lr@{}}
\toprule
Placas                               & \multicolumn{1}{l}{Resultado} \\ \midrule
Total de Imagens                     & 100                           \\
Imagens onde a placa foi encontrada  & 10                            \\
Imagens onde a placa foi reconhecida & 1                            
\end{tabular}
\end{table}

\section{Extração da placa}
\label{sec:extracao_da_placa_resultados}


O algoritmo de extração da placa é limitado com relação a distância que o carro está na imagem. Na etapa de abertura e fechamento morfológicos, que tem o objetivo de detectar as áreas de placas candidatas, o tamanho do elemento estruturante influencia diretamente a qualidade da extração. Se utilizado um elemento estruturante muito grande, é possível que na etapa de abertura, a placa seja erodida junto com o ruído. Se utilizado um elemento estruturante muito pequeno, muitas regiões sobram, diminuindo a performance do reconhecimento. 

Colocar aqui imagem de exeplo mostrando a placa sendo destruida

\section{Reconhecimento dos caracteres}
\label{sec:reconhecimento_dos_caracteres_resultados}

Inicialmente foi utilizado o \emph{software Tesseract} para fazer o reconhecimento dos caracteres, mas os resultados utilizando a ferramenta pura não foram bons. Mesmo limitando os caracteres possíveis, informando ao \emph{software} se o caractere seria um número ou uma letra, ainda havia dificuldade no reconhecimento. Mesmo em caracteres que não são tão semelhantes, como o 5 e o 6, havia dificuldade no reconhecimento.

Com a substituição do reconhecedor de caracteres \emph{Tesseract} para o implementado com o algoritmo \emph{K-Nearest Neighbors} os resultados foram bem melhores. Até mesmo em imagens de placas inclinadas, como é o caso da Figura~\ref{fig:plate_torta_result}, foi possível identificar os caracteres.

\begin{figure}[H]
	\centering
	\includegraphics[width=88mm]{a42fill_binary_results.jpg}
	\caption{Placa inclinada que o sistema foi capaz de identificar}
	\label{fig:plate_torta_result}
\end{figure}

Colocar aqui os resultados e quais caracteres foram confundidos e em quais condições

\section{Performance em sistema embarcado}
\label{sec:performance_resultados}

Dado o fato de a solução desenvolvida ser dividida em etapas, é possível avaliar quais os passos levam mais tempo para serem executados, que podem ser considerados os gargalos da aplicação. Em uma visão mais ampla, fica claro que a etapa de extração da placa é a que necessita de mais recursos computacionais. Isso se dá principalmente por ter que processar uma imagem maior que a imagem da placa. O tempo utilizado para processar a placa extraída é, relativamente, bem menor.

Para melhor analisar os resultados comparativos, foram calculadas as médias de tempo de execução de cada etapa em múltiplas imagens e calculados seus valores percentuais em relação ao tempo total de processamento de uma placa. Esses valores independem do \emph{hardware} utilizado e do tamanho de imagem utilizado, permitindo uma visão diferente dos resultados. Os resultados comparativos podem ser vistos na Tabela~\ref{tab:resultados_relativos}.

\begin{table}[H]
\centering
\caption{Resultados relativos da fase de extração}
\label{tab:resultados_relativos}
\begin{tabular}{@{}lr@{}}
\toprule
Etapa                        & Tempo Percentual \\ \midrule
Conversão para tons de cinza & 10\%             \\
Filtro Bilateral             & 10\%             \\
Equalização do Histograma    & 10\%             \\
Limiarização                 & 10\%             \\
Detecção de Bordas           & 10\%             \\
Dilatação                    & 10\%             \\
Preenchimento de espaços     & 10\%             \\
Abertura                     & 10\%             \\
Fechamento                   & 10\%             \\
Extração das Regiões         & 10\%             \\ \bottomrule
\end{tabular}
\end{table}

A primeira implementação do algoritmo teve maus resultados com relação a performance. 

Falar sobre os resultados aqui. 

Para tentar aprimorar a velocidade do processamento em imagens consecutivas foi aplicado o uso de múltiplos processos, com o objetivo de utilizar todos os núcleos do processador do \emph{Raspberry Pi}.
